\documentclass[11pt, a4paper]{article}

\usepackage[T1]{fontenc}
\usepackage[utf8]{inputenc}
\usepackage{tgpagella}
\usepackage{tgcursor}
\usepackage{microtype}
\usepackage[spanish, es-noshorthands]{babel}
\usepackage[margin=2.5cm]{geometry}
\usepackage{fancyhdr}
\usepackage[most]{tcolorbox}
\usepackage{listings}

\lstset{
    language=Python,
    basicstyle=\ttfamily\small,
    backgroundcolor=\color{gray!5},
    frame=single,
    rulecolor=\color{gray!40},
    keywordstyle=\color{blue}\bfseries,
    stringstyle=\color{red!70!black},
    commentstyle=\color{green!60!black},
    showstringspaces=false,
    breaklines=true,          % Permite salto de linea automatico
    breakatwhitespace=true,   % Salta solo en espacios en blanco
    columns=flexible          % Ajusta el espaciado de los caracteres
}

\pagestyle{fancy}
\fancyhf{}
\setlength{\headheight}{16pt}
\lhead{\textbf{UCB Sede Tarija} | Ing. Mecatrónica}
\chead{Código: Jacobiano en Python}
\rhead{Opcional W08-S01}

\definecolor{ucbblue}{RGB}{0, 51, 102}

\begin{document}

\begin{center}
    \Large\textbf{\textcolor{ucbblue}{Verificación Computacional: Jacobiano Analítico 2R}}\\[0.2cm]
\end{center}

Este script correlaciona el modelo matemático simbólico con la evaluación numérica, mapeo de velocidades y cálculo de rango (para validación de singularidades)[cite: 11]. 
\textit{Dependencias requeridas: Python 3 y Numpy.}

\begin{lstlisting}
import numpy as np

def jacobian_2r(q, a1=1.0, a2=1.0):
    """ 
    Retorna el sub-jacobiano de traslacion 
    2x2 para brazo 2R 
    """
    q1, q2 = q
    
    s1 = np.sin(q1)
    c1 = np.cos(q1)
    s12 = np.sin(q1 + q2)
    c12 = np.cos(q1 + q2)
    
    # Ensamblaje del modelo simbolico derivado en clase
    return np.array([
        [-a1*s1 - a2*s12, -a2*s12],
        [ a1*c1 + a2*c12,  a2*c12]
    ], dtype=float)

# ==========================================
# 1. PRUEBA DE MAPEO DE VELOCIDAD
# ==========================================

# q1=0, q2=90 grados
q = np.array([0.0, np.pi/2])  

# qdot: velocidades articulares comandadas
qd = np.array([1.0, -1.0])    

J = jacobian_2r(q)
pd = J @ qd                   # Map: J * q_dot

print("--- Configuracion Regular ---")
print("Matriz J =\n", np.round(J, 4))
print("Vel. Cartesiana (p_dot) =", 
      np.round(pd, 4))
print("Rango(J) =", np.linalg.matrix_rank(J))

# ==========================================
# 2. PRUEBA DE SINGULARIDAD (Brazo estirado)
# ==========================================

# q2 = 0 es configuracion singular analitica
q_sing = np.array([0.0, 0.0]) 
J_sing = jacobian_2r(q_sing)

print("\n--- Configuracion Singular ---")
print("Matriz J_sing =\n", np.round(J_sing, 4))
print("Rango(J_sing) =", np.linalg.matrix_rank(J_sing))
\end{lstlisting}

\begin{tcolorbox}[colback=white, colframe=ucbblue, title=\textbf{Salida Esperada y Correspondencia}]
\textbf{Configuración Regular:} La salida de \texttt{p\_dot} será \texttt{[0. 1.]}, confirmando la cancelación de la traslación en X (rango = 2, control cartesiano completo). \\
\textbf{Configuración Singular:} La matriz tendrá columnas dependientes \texttt{[[0. 0.], [2. 1.]]}, y el rango caerá a \texttt{1}, comprobando matemáticamente que el robot perdió movilidad instantánea independiente en una dirección[cite: 11].
\end{tcolorbox}

\end{document}
